% Generated draft: requires mathematical and source review.
\documentclass[11pt]{article}
\usepackage[margin=25mm]{geometry}
\usepackage{fontspec}
\setmainfont{Latin Modern Roman}
\usepackage{amsmath,amssymb,mathtools,unicode-math}
\setmathfont{Latin Modern Math}
\usepackage{enumitem,needspace,xurl,hyperref}
\setlength{\parindent}{0pt}
\setlength{\parskip}{5pt}
\setlength{\emergencystretch}{3em}
\newcommand{\N}{\mathbb{N}}
\newcommand{\Z}{\mathbb{Z}}
\newcommand{\Q}{\mathbb{Q}}
\newcommand{\R}{\mathbb{R}}
\newcommand{\C}{\mathbb{C}}
\newcommand{\F}{\mathbb{F}}
\newcommand{\A}{\mathbb{A}}
\newcommand{\PP}{\mathbb{P}}
\newcommand{\E}{\mathbb{E}}
\newcommand{\eps}{\varepsilon}
\newcommand{\dd}{\,\mathrm{d}}
\begin{document}
\section*{UnsolvedMath ID: 30001554}
\textbf{Source problem code:} catalogue entry
\subsection{Definitions and mathematical statement}
Let \(A\) be an alphabet, and let \(A^*\) denote the monoid of finite words over \(A\), with concatenation and empty word \(\varepsilon\). For a finite word \(w\), let \(|w|\) denote its length. A morphic involution is a map \(\theta:A^*\to A^*\) satisfying
\[
\theta(\varepsilon)=\varepsilon,\qquad
\theta(uv)=\theta(u)\theta(v)\quad(u,v\in A^*),\qquad
\theta^2=\operatorname{id}_{A^*}.
\]
The supplied question asks whether, for every such \(\theta\) and every \(A\)-word \(w\),
\[
|w|\geq 3\tau_\theta(w)
\quad\Longrightarrow\quad
\tau_\theta(w)=\pi_\theta^{\mathrm{alt}}(w).
\]
\textbf{Unresolved definitions and scope.} The supplied record does not define either \(\tau_\theta(w)\) or \(\pi_\theta^{\mathrm{alt}}(w)\), including their domains, boundary conventions, or the precise meaning of an alternating involution period. Consequently, the displayed implication preserves the source question but is not a fully specified mathematical conjecture. The title mentions unbordered factors, but it does not determine a definition of \(\tau_\theta\). No finiteness restriction on \(A\), restriction on fixed letters of \(\theta\), or explicit treatment of the empty word is supplied. These conventions must be recovered from the cited source before the universal assertion can be made precise.
\subsection{Short English statement}
The question asks whether a word whose length is at least three times \(\tau_\theta(w)\) must satisfy \(\tau_\theta(w)=\pi_\theta^{\mathrm{alt}}(w)\), for a morphic involution \(\theta\). The two numerical quantities are undefined in the supplied record, so their intended definitions and boundary conventions require human review.
\subsection{Sources}
\begin{itemize}
\item Local UnsolvedMath registry: supplied record with permanent ID \(30001554\), source problem code \texttt{catalogue entry}, title ``Unbordered Factors and Alternating Involution Periods,'' in the Oberwolfach Reports Open Problems collection. No individual proposer is supplied; the proposed year is recorded as 2010.
\item Supplied source citation: \emph{Mini-Workshop: Combinatorics on Words (2011)}, Oberwolfach Reports, pp.\ 2195--2244, DOI 10.4171/owr/2010/37. Original workshop report: \url{https://doi.org/10.4171/owr/2010/37}. The background describes these reports as documenting mathematical workshops and their open-problem discussions.
\item Supplied later thesis exposition: \url{https://www2.informatik.uni-stuttgart.de/bibliothek/ftp/medoc_restrict.ustuttgart_fi/DIP-3095/DIP-3095.pdf}. No thesis author, title, or relevant page is supplied.
\item Literature assessment dated exactly \texttt{2026-08-21}: the supplied curation labels the problem open and reports that the proposed extremal relation between theta-periods and alternating theta-periods remains conjectural in later expositions, with no proof or counterexample found. This is supplied dated literature triage derived from the OWR curation, not an independently verified assessment. The record lists the original workshop report and later thesis exposition as sources checked; time-sensitive status and cited claims require independent verification.
\item The supplied statement status is \texttt{corrected\_verified}; its accompanying verification note says only that TeX notation was enclosed in mathematical delimiters without changing the mathematical content. This records a formatting correction and does not establish mathematical verification.
\end{itemize}
\end{document}
